\section{周王室之二：当中心无法再调停各方}

\begin{event}{公元前841年}{厉王出奔与共和行政}{可靠纪年}{镐京}\eventanchor{WZ-0841-GONGHE}
\term{这一年发生了什么}：厉王出奔，周、召二公共同执政，传统纪年称为“共和”。这也是传统中国连续纪年的确切起点：从这一年开始，历史叙事第一次拥有一根可以逐年往后数的稳定标尺。

\term{事情为什么走到这里}：\sourcebook{国语}等材料将矛盾放在王室资源政策与“防民之口”的冲突上。无论后世叙事是否放大了单一政策，核心问题很清楚：王室既需要资源维持军政与赏赐，又失去了让贵族和国人接受这种分配的政治能力。

\term{当事人的选择与代价}：厉王若继续以压制回应批评，短期或许能保住政策，长期却会把可协商的不满变成共同反抗。周、召共同执政保住了王位继承的形式，却也公开显示王权已不能独立维持秩序。

\term{后来改变了什么}：共和并未让西周立即崩溃，却留下一个更清楚的信号：王室与其核心支持者之间的信任已被削薄。它是\eventid{WZ-0771-HAOJING}之前最醒目的内部裂缝。
\end{event}

\begin{textwindow}[“防民之口，甚于防川”]
《国语·周语上》记召公谏厉王：“防民之口，甚于防川。川壅而溃，伤人必多；民亦如之。”这段话之所以流传，不只是因为它像一条现代言论格言，而是因为它把西周晚期的政治难题说得很直白：堵住表达未必消除不满，只会让不满失去可见、可谈、可调停的出口。
\end{textwindow}

\begin{debate}[“共和”是谁在执政]
传统解释多称周公、召公共同行政；也有材料和研究提出“共伯和”说。这里不把一种说法写成最终答案。对本书的叙事而言，更重要的是：王室在厉王出奔后需要一个非正常的政治安排来维持连续性。
\end{debate}

\begin{event}{公元前771年}{镐京陷落，幽王死}{可靠纪年}{镐京}\eventanchor{WZ-0771-HAOJING}
\term{这一年发生了什么}：幽王死于申侯、犬戎等势力攻入镐京的危机中。次年，平王迁都洛邑。西周作为以关中为核心的王室秩序至此结束。

\term{事情为什么走到这里}：传统叙事聚焦幽王废申后与太子宜臼、改立伯服。继承冲突确实把王室内部问题公开化；但它能引发如此严重的后果，也因为西部边防、王室军力、封国网络和内部信任早已同时松动。

\term{当事人的选择与代价}：继承从来不是私人的家庭决定。王后、太子、外戚、申侯与边疆势力各自判断风险后，原本的宗法秩序变成了动员不同联盟反对王室的工具。

\term{后来改变了什么}：平王东迁后的周仍是“天子”，但已从资源和军事中心退为礼仪与名分中心。春秋诸侯争霸、会盟和“尊王攘夷”的语言，都建立在这个中心还在、力量却已不足的矛盾上。
\end{event}

\begin{causalchain}[从共和到东迁]
\eventid{WZ-0841-GONGHE} 显示王室不能稳住内部；\eventid{WZ-0771-HAOJING} 则让继承危机、诸侯联盟和边防失效在同一时刻爆发。东迁不是一场搬家，而是统治能力的重心转移：周王室保住了名分，却失去了支配名分的物质底座。
\end{causalchain}
